% GENERATED FILE. Edit canonical lessons, then run npm run generate:books.
% canonical-source-hash: fnv1a64:9d2fabac10dea04a
% canonical-lessons: RU-C519-zanaves, RU-C519-antrakt, RU-C519-aplodirovat, RU-C519-premyera, RU-C519-klanyatsya

\chapter{Curtain Up}
\label{ch:ru-curtain-up}
\hlchaptermodality{519}

\begin{chapteropening}
\textbf{By the end of this chapter:} \emph{I can talk about an evening at the theatre.}

Complete the last lesson of chapter 519: I can talk about an evening at the theatre.
\end{chapteropening}

\section[\texorpdfstring{\ru{занавес} (zánaves)}{занавес (zánaves)}]{\ru{занавес} (zánaves) — a stage curtain}

\label{lesson:RU-C519-zanaves}



\begin{quote}
Before the new one: say the Russian for to rehearse, then the Russian for a tour.
\end{quote}

\subsection*{You'll want to know: \ru{занавес}}
\textbf{\ru{занавес}} (\emph{zánaves}) — \textquotedblleft{}a stage curtain\textquotedblright{}.

\begin{quote}
\ru{Занавес} \ru{поднялся}, \ru{и} \ru{спектакль} \ru{начался}.
\end{quote}

\begin{quote}
\emph{Zánaves podnyálsya, i spektákl nachalsyá.}
\end{quote}

\begin{quote}
The curtain went up and the play began.
\end{quote}

Masculine. A window curtain is \emph{\ru{штора}} or \emph{\ru{занавеска}}.

\begin{cousinweb}
\emph{\ru{занавес}} is derived from \emph{\ru{навес}} 'an awning' (EtymDB), built on the root of \emph{\ru{вешать}} 'to hang'.
\end{cousinweb}

\begin{grammarlens}[title={What you've built}]
One more word for telling, explaining and describing at B1.
\end{grammarlens}

\subsection*{Your turn}
\begin{itemize}
\raggedright
  \item \emph{Say it:} \emph{zánaves}
  \item \emph{Say it:} \emph{zánaves}, once more
  \item \emph{Say it:} a sentence of your own with \emph{zánaves}
  \item \emph{Recall:} say the Russian for to rehearse, then the Russian for a tour, then say \emph{zánaves} again
\end{itemize}

\begin{tcolorbox}[breakable,colback=teal!4,colframe=teal!35!black,title={Before you move on}]
What does \ru{занавес} mean? (\textquotedblleft{}A stage curtain\textquotedblright{}.) Say it once more.
\end{tcolorbox}
\section[\texorpdfstring{\ru{антракт} (antrákt)}{антракт (antrákt)}]{\ru{антракт} (antrákt) — an interval}

\label{lesson:RU-C519-antrakt}



\begin{quote}
Before the new one: say the Russian for a tour, then the Russian for a stage curtain.
\end{quote}

\subsection*{You'll want to know: \ru{антракт}}
\textbf{\ru{антракт}} (\emph{antrákt}) — \textquotedblleft{}an interval\textquotedblright{}.

\begin{quote}
\ru{В} \ru{антракте} \ru{мы} \ru{пили} \ru{кофе} \ru{в} \ru{буфете}.
\end{quote}

\begin{quote}
\emph{V antrákte my píli kófe v buféte.}
\end{quote}

\begin{quote}
In the interval we had coffee at the bar.
\end{quote}

Masculine.

\begin{cousinweb}
\emph{\ru{антракт}} was borrowed from French \emph{entr'acte} 'the time between the acts' (EtymDB).
\end{cousinweb}

\begin{grammarlens}[title={What you've built}]
One more word for telling, explaining and describing at B1.
\end{grammarlens}

\subsection*{Your turn}
\begin{itemize}
\raggedright
  \item \emph{Say it:} \emph{antrákt}
  \item \emph{Say it:} \emph{antrákt}, once more
  \item \emph{Say it:} a sentence of your own with \emph{antrákt}
  \item \emph{Recall:} say the Russian for a tour, then the Russian for a stage curtain, then say \emph{antrákt} again
\end{itemize}

\begin{tcolorbox}[breakable,colback=teal!4,colframe=teal!35!black,title={Before you move on}]
What does \ru{антракт} mean? (\textquotedblleft{}An interval\textquotedblright{}.) Say it once more.
\end{tcolorbox}
\section[\texorpdfstring{\ru{аплодировать} (aplodírovat)}{аплодировать (aplodírovat)}]{\ru{аплодировать} (aplodírovat) — to applaud}

\label{lesson:RU-C519-aplodirovat}



\begin{quote}
Before the new one: say the Russian for a stage curtain, then the Russian for an interval.
\end{quote}

\subsection*{You'll want to know: \ru{аплодировать}}
\textbf{\ru{аплодировать}} (\emph{aplodírovat}) — \textquotedblleft{}to applaud\textquotedblright{}.

\begin{quote}
\ru{Публика} \ru{долго} \ru{аплодировала} \ru{актёрам}.
\end{quote}

\begin{quote}
\emph{Públika dólgo aplodírovala aktyóram.}
\end{quote}

\begin{quote}
The audience applauded the actors for a long time.
\end{quote}

Imperfective; present \ru{аплодирую}. Takes the dative. The noun is \emph{\ru{аплодисменты}} — applause.

\begin{cousinweb}
\emph{\ru{аплодировать}} is a learned borrowing, usually traced through French \emph{applaudir} to Latin \emph{applaudere} 'to applaud' (Lewis \& Short); EtymDB lists the verb without a source.
\end{cousinweb}

\begin{grammarlens}[title={What you've built}]
One more word for telling, explaining and describing at B1.
\end{grammarlens}

\subsection*{Your turn}
\begin{itemize}
\raggedright
  \item \emph{Say it:} \emph{aplodírovat}
  \item \emph{Say it:} \emph{aplodírovat}, once more
  \item \emph{Say it:} a sentence of your own with \emph{aplodírovat}
  \item \emph{Recall:} say the Russian for a stage curtain, then the Russian for an interval, then say \emph{aplodírovat} again
\end{itemize}

\begin{tcolorbox}[breakable,colback=teal!4,colframe=teal!35!black,title={Before you move on}]
What does \ru{аплодировать} mean? (\textquotedblleft{}To applaud\textquotedblright{}.) Say it once more.
\end{tcolorbox}
\section[\texorpdfstring{\ru{премьера} (premyéra)}{премьера (premyéra)}]{\ru{премьера} (premyéra) — a premiere}

\label{lesson:RU-C519-premyera}



\begin{quote}
Before the new one: say the Russian for an interval, then the Russian for to applaud.
\end{quote}

\subsection*{You'll want to know: \ru{премьера}}
\textbf{\ru{премьера}} (\emph{premyéra}) — \textquotedblleft{}a premiere\textquotedblright{}.

\begin{quote}
\ru{Премьера} \ru{спектакля} \ru{будет} \ru{в} \ru{субботу}.
\end{quote}

\begin{quote}
\emph{Premyéra spektáklya búdet v subbótu.}
\end{quote}

\begin{quote}
The premiere of the play is on Saturday.
\end{quote}

Feminine.

\begin{cousinweb}
\emph{\ru{премьера}} was borrowed from French \emph{première} 'first' (EtymDB).
\end{cousinweb}

\begin{grammarlens}[title={What you've built}]
One more word for telling, explaining and describing at B1.
\end{grammarlens}

\subsection*{Your turn}
\begin{itemize}
\raggedright
  \item \emph{Say it:} \emph{premyéra}
  \item \emph{Say it:} \emph{premyéra}, once more
  \item \emph{Say it:} a sentence of your own with \emph{premyéra}
  \item \emph{Recall:} say the Russian for an interval, then the Russian for to applaud, then say \emph{premyéra} again
\end{itemize}

\begin{tcolorbox}[breakable,colback=teal!4,colframe=teal!35!black,title={Before you move on}]
What does \ru{премьера} mean? (\textquotedblleft{}A premiere\textquotedblright{}.) Say it once more.
\end{tcolorbox}
\section[\texorpdfstring{\ru{кланяться} (klányatsya)}{кланяться (klányatsya)}]{\ru{кланяться} (klányatsya) — to bow}

\label{lesson:RU-C519-klanyatsya}



\begin{quote}
Before the new one: say the Russian for to applaud, then the Russian for a premiere.
\end{quote}

\subsection*{You'll want to know: \ru{кланяться}}
\textbf{\ru{кланяться}} (\emph{klányatsya}) — \textquotedblleft{}to bow\textquotedblright{}.

\begin{quote}
\ru{В} \ru{конце} \ru{спектакля} \ru{актёры} \ru{кланяются}.
\end{quote}

\begin{quote}
\emph{V kontsé spektáklya aktyóry klányayutsya.}
\end{quote}

\begin{quote}
At the end of the play the actors take a bow.
\end{quote}

Imperfective; perfective \ru{поклониться}.

\begin{cousinweb}
EtymDB lists \emph{\ru{кланяться}} without a source; it is usually linked to the root of \emph{\ru{клонить}} 'to bend, to tilt'.
\end{cousinweb}

\begin{grammarlens}[title={What you've built}]
One more word for telling, explaining and describing at B1.
\end{grammarlens}

\subsection*{Your turn}
\begin{itemize}
\raggedright
  \item \emph{Say it:} \emph{klányatsya}
  \item \emph{Say it:} \emph{klányatsya}, once more
  \item \emph{Say it:} a sentence of your own with \emph{klányatsya}
  \item \emph{Recall:} say the Russian for to applaud, then the Russian for a premiere, then say \emph{klányatsya} again
\end{itemize}

\begin{tcolorbox}[breakable,colback=teal!4,colframe=teal!35!black,title={Before you move on}]
What does \ru{кланяться} mean? (\textquotedblleft{}To bow\textquotedblright{}.) Say it once more.
\end{tcolorbox}
